\documentclass{article}

\usepackage{fullpage}
\usepackage{parskip}
\usepackage{setspace}
\usepackage{amsmath}
\usepackage{amsfonts}
\usepackage{amsthm}
\usepackage{hyperref}
\usepackage{xcolor}

\title{Responses to Reviewers}
\author{}
\date{}


\begin{document}

\maketitle

We would like to express our gratitude to the editor, associate editor, and reviewers for facilitating the review process, and for their useful suggestions and remarks on the manuscript.

This document will address the specific concerns raised.

\section{Responses to the Reviewers}

\subsection{Comments from Reviewer 1}

\begin{itemize}

\item Reviewer 1 wrote:
\begin{quote}\it
``Motivation: The authors use the healthcare system as a motivating example, illustrating that a patient’s priority can upgrade, downgrade, or skip-grade. For medical urgency, priority upgrades are easy to understand. However, what is the meaning and context of `skip-grades'? What is the precise definition of priority? The authors also state that medical intervention can improve a patient's outcome and that the patient’s priority can downgrade. This implies that the patient receives the service (treatment) and therefore is no longer waiting in the queue.''
\end{quote}

\textcolor{red}{To address...}

\item Reviewer 1 wrote:
\begin{quote}\it
``The authors use two separate Markov chains to find state probabilities and sojourn times. Can the authors attempt to analyze the model theoretically and provide corresponding theoretical results? If not, perhaps they could explain the difficulty of the problems. In addition, the authors should explain Equation (1) and Equation (5) more clearly to readers. To use Markov chains, the authors present bounded approximations. Can the authors explicitly analyze this bounded approximation method? For example, could they theoretically determine the value of $b$ and prove the system's stability?''
\end{quote}

\textcolor{red}{To address...}


\item Reviewer 1 wrote:
\begin{quote}\it
``Most of the results in this paper are obtained from the numerical studies. However, the author does not explain the source of the parameters in detail, that is, why they were chosen. Also, I was confused for prioritised customers have a faster service rate. The authors provide a motivating case study in section 1.1, ideally the parameters should come from this case study.''
\end{quote}

\textcolor{red}{To address...}

\item Reviewer 1 wrote:
\begin{quote}\it
``Page 1, line 26: ``There are a number of situations in which a customer's priority in a queue might change during their time queueing...'' should be ``There are several situations in which a customer's priority in a queue might change while they wait...''''
\end{quote}

\textcolor{red}{To address...}

\item Reviewer 1 wrote:
\begin{quote}\it
``Page 5, line 46: ``and extra C-events would include include...''''
\end{quote}

\textcolor{red}{Thank you, this has been fixed.}

\item Reviewer 1 wrote:
\begin{quote}\it
``Page 6, line 42: ``According to the observed data referrals roughly follow a Poisson dis- tribution with rate 0.463 a day. Service rate data is estimated to be around 0.5 endoscopy procedures a day.'' What it the meaning of 0.463 and 0.5, why not integers? Also, the authors assume that all referrals are of the most urgent, with customers able to be downgraded and then upgraded during their waiting time may not be realistic.''
\end{quote}

\textcolor{red}{To address...}

\item Reviewer 1 wrote:
\begin{quote}\it
``Page 11, line 50: ``Section 4.2)'' should be ``Section 4.2''.''
\end{quote}

\textcolor{red}{To address...}

\item Reviewer 1 wrote:
\begin{quote}\it
``Page 12, Figure 8: what is the meaning of ``Class 0'', ``Class 1'' and ``Combined''.''
\end{quote}

\textcolor{red}{To address...}

\item Reviewer 1 wrote:
\begin{quote}\it
``Page 19, line 36: ``...demonstrating some important performance..'' should be ``...demonstrating that some important performance..''.''
\end{quote}

\textcolor{red}{To address...}

\end{itemize}


\subsection{Comments from Reviewer 2}

\begin{itemize}

\item Reviewer 2 wrote:
\begin{quote}\it
``The article is studying an interesting topic -- stochastic priority is not well explored in the literature, despite its wide applications (e.g., healthcare). The insights of the paper well complements the existing literature on dynamic priority (or accumulating priority queue).''
\end{quote}

\textcolor{red}{To address...}

\item Reviewer 2 wrote:
\begin{quote}\it
``I find Section 4.6, the study of the impact of stochastic priority on variance of sojourn time, particularly interesting. But the current exploration along this direction is rather limited -- the paper only studied one problem instances with two priority classes. Can the authors provide more comparisons between stochastic priority and non-stochastic priority under various parameters (e.g., other $\theta_{ij}$, and more priority classes) as robustness test? Also, in addition to variance of waiting time, maybe the 95\% percentile waiting time is more accessible to readers?''
\end{quote}

\textcolor{red}{To address...}

\item Reviewer 2 wrote:
\begin{quote}\it
``The simulation in Section 3 is not very interesting. Can these materials regarding setup of the simulation be deferred to appendix? The simulation setup is quite straightforward.''
\end{quote}

\textcolor{red}{To address...}

\item Reviewer 2 wrote:
\begin{quote}\it
``In conclusion, the paper is studying an interesting topic, providing a good complement to the existing literature. The case study and the insight is interesting and intriguing. I thus recommend a minor revision.''
\end{quote}

\textcolor{red}{To address...}

\end{itemize}

\end{document}